\section{Introduction}\label{sec:introduction}
Let $X$ be a smooth connected projective variety over $\C$ and put $L=-K_X=\det T_X$. Anticanonical nonvanishing asks whether $H^0(X,mL)\ne0$ for some integer $m>0$. We assume that $L$ is smoothly semipositive: it has a smooth Hermitian metric with nonnegative Chern curvature. Cohomological constructions can produce sections with a fixed twist, which must then be removed. In the presence of a group action there is a second obstruction: the resulting section must have the required character. Our principal theorem addresses these two obstructions together.

Let an algebraic torus $T$ act on $X$. A $T$-linearization of a line bundle is a lift of this action to its total space, linear on each fiber. The tangent action gives the natural linearization of $L$, but tensoring that action by a character of $T$ gives another one. We write $P$ for $L$ equipped with any such linearization, and use additive notation for tensor products. Thus $H^0(X,mP)^T$ specifies both a degree and a character of an anticanonical section.

The input to our principal theorem is a sequence of sections of $M+iP$ with a fixed error $M$. The assumption on the error is that $-M$ is pseudoeffective: its numerical class lies in the closure of the effective cone.

\begin{theorem}[Invariant transfer]\label{inj:transfer}
Let $X$ be smooth connected projective, let $L=-K_X$ carry a smooth
semipositive metric, and let an algebraic torus $T$ act on $X$.
Let $P,M$ be linearized line bundles such that $P=L$ as ordinary
line bundles and $-M$ is pseudoeffective.  The linearization of $P$
may differ from the natural one on $L$.  Suppose that for an
unbounded $I\subset\mathbb Z_{>0}$ there are nonzero invariant
sections
\[
 s_i\in H^0(X,M+iP)^T\qquad(i\in I).
\]
Then $H^0(X,mP)^T\ne0$ for some integer $m>0$.
\end{theorem}


The theorem removes the fixed error without losing the specified character. It assumes neither vanishing of $H^1(X,\mathcal O_X)$ nor a condition on the Euler characteristic. Taking the trivial torus gives ordinary fixed-error conversion. With the natural linearization, invariant cohomology supplies an important source of the required twisted sections. In Section~\ref{sec:cohomological-input} we deduce that
\[
 \chi(X,\mathcal O_X)\ne0
 \quad\Longrightarrow\quad
 H^0(X,-mK_X)^T\ne0\text{ for some }m>0.
\]
The integer may depend on $X$. The distinction between natural and shifted linearizations remains essential in this deduction: the invariant anticanonical index is applied only with the natural one.

Combining this natural invariant nonvanishing with the compact-monodromy structure theorem of \cite[Theorem \textnormal{(Finite cover with compact torus monodromy)}]{CompanionH} removes the Euler-characteristic premise from the ordinary conclusion. The structure theorem separates the universal cover into a Euclidean factor, compact factors with parallel canonical frames, and a compact factor with no positive-degree holomorphic forms. After a finite cover, the remaining action on the last factor lies in a compact torus. Its invariant section therefore descends, and a finite norm returns a section to $X$. Section~\ref{sec:descent} proves the following consequence and constructs two additional descent maps, for cohomology and for twisted forms.

\begin{corollary}\label{thm:ordinary}
For every smooth connected projective complex $X$ with smoothly semipositive $-K_X$, some positive multiple $-mK_X$ has a nonzero section.
\end{corollary}

\subsection{Why a fibration removes the error}
Choose nonzero invariant input sections $s_i$. Ratios of products having the same linearized degree are invariant rational functions. Their common field determines a rational map from $X$ to a projective base. Along its generic fiber, the divisors of the $s_i$ vary affinely with $i$. Unboundedness forces their slope to be effective, so a quotient $s_j/s_i$ supplies a relative anticanonical section.

The main difficulty is to extend this relative section across the boundary. We work with a birational morphism $\mu:W\to X$ and a normal equidimensional model $f_W:W\to Y$, and put $P_W=\mu^*P$. Every vertical prime of $W$ then lies over a prime of $Y$. Subtracting the smallest normalized order over each base prime gives a regular section of $kP_W-f_W^*(kB)$, where $B$ is a rational divisor on $Y$ and $kB$ is integral. On a smooth resolution $\pi:V\to X$ through $W$, put $E=K_{V/X}$. The same ratio field shows that the invariant sections of $E+jk\pi^*P$ on the generic fiber of $V\to Y$ form a one-dimensional space over $\C(Y)$ for every $j\ge0$, including zero.

Two metrics now play different roles. Smooth adjoint fiber integrals in increasing degrees converge, after taking roots, to a fiberwise maximum. This gives a bounded semipositive metric on $kB$. At degree zero, pseudoeffectivity of $-M$ supplies a different metric with positive curvature in the base directions. Its fiber integral extends to $-K_Y+D_*$, where $D_*\ge0$ pulls back to an exceptional divisor over $X$. Its multiplier ideal $\mathcal J$ can be nontrivial. What matters is the inclusion $\mathcal O_Y(-D_*)\subseteq\mathcal J$ and the fact that tensoring by the bounded metric on $jkB$ leaves $\mathcal J$ unchanged.

Nadel vanishing therefore turns the section dimension of
$\mathcal O_Y(D_*+jkB)\otimes\mathcal J$ into a polynomial in $j$, including at zero. Its value there is positive, so it cannot vanish at every positive integer. A resulting section lifts by the relative section, and the exceptional correction disappears on $X$. Section~\ref{sec:models} prepares the model and the maximum metric; Section~\ref{inj:character-transfer} gives the complete transfer proof.

\subsection{A separate harmonic mechanism}
There is another way to detect the zero exponent, using restriction of cohomology to fibers. Its hypotheses and purpose differ from character-sensitive transfer. Let $f:V\to Y$ be a surjective morphism between smooth projective varieties, let $D_V$ have a smooth semipositive metric of curvature $\theta$, and let $Q$ have a continuous semipositive metric on $Y$. Suppose that on a nonempty ordinary open set where $f$ is a submersion, every null direction of $\theta$ is vertical. We prove in Theorem~\ref{har:thm:kernel} that
\[
 H^p\bigl(Y,R^qf_*(K_V+D_V)\otimes Q^k\bigr)=0
 \qquad(p>0,\ q,k\ge0).
\]
The proof approximates only the continuous metric pulled back from $Y$. A weak limit of harmonic representatives has all base canonical covectors as wedge factors, and its relative part detects a nonzero cohomology class on a general fiber. We apply decomposition to the smooth coefficient $D_V$ and then tensor the resulting derived isomorphism by $Q^k$. Thus the argument uses no decomposition theorem for a merely continuous coefficient upstairs.

We work through the natural-linearization application completely, then explain alternative constructions at the point where their geometric input changes. A fixed invariant system, null-pole functions, and determinant-semigroup ratios give different ways to choose the base. An ordinary curvature-null field leads instead to a reduced flat alteration, a line of multiplication maps, and a norm back to the original function field. Local strictness with an exact direct-image ideal and stabilization of actual direct-image lines provide two further vanishing mechanisms. Each construction is carried through to its section conclusion; the common analytic conclusions are reused only after their hypotheses have been verified.

\subsection{Context and analytic predecessors}
Anticanonical nonvanishing is motivated by the questions recorded by Yau \cite[Problem~75]{Yau1994}. Our conclusions concern a positive power, whose exponent may vary with the variety. Smooth semipositivity provides curvature identities unavailable from numerical nefness alone. Demailly--Peternell--Schneider established the universal-cover splitting and finite-cover Albanese fibration for a Hermitian semipositive anticanonical bundle \cite[Structure Theorem]{DPS1996}; Campana--Demailly--Peternell identified the remaining compact factors as rationally connected \cite[Theorem~1.4]{CDP2015}. The finite cover is a fiber bundle, so descent must still account for monodromy. We use the precise compact-torus structure and finite \'etale norm in \cite{CompanionH}, and construct the cohomology-valued and form-valued descent maps here.

In the broader nef setting, Lazi\'c--Matsumura--Peternell--Tsakanikas--Xie proved numerical effectivity of $-(K_X+\Delta)$ for projective log canonical threefold pairs with nef $-(K_X+\Delta)$, assuming that $X$ is $\Q$-factorial or has rational singularities \cite[Theorem~A]{LMPTX2023}. Numerical effectivity means that the numerical class is represented by an effective $\R$-divisor; it does not by itself provide a section of the specified line bundle.

M\"uller proves nonvanishing for projective klt pairs with nef log anticanonical divisor whose restriction to a general fiber of the maximal rationally connected fibration is semiample, and for every projective klt threefold pair with nef log anticanonical divisor \cite[Theorem~A and Corollary~B]{Muller2025}. His equivariant theorem gives natural invariant sections for semiample log anticanonical divisors on projective sub-log canonical pairs under commutative linear algebraic groups \cite[Theorem~C]{Muller2025}. Here we work with smooth projective varieties and algebraic tori. Smooth semipositivity and a nonzero Euler characteristic give natural invariant sections, while Theorem~\ref{inj:transfer} allows shifted linearizations with fixed-error input. The compact-monodromy descent then gives ordinary nonvanishing without the Euler-characteristic assumption.

The independent compact-complex index theorem, with its actual value at zero, is supplied by \cite[Invariant anticanonical index theorem]{CompanionI}. The passage from cohomology to twisted forms uses smooth semipositive hard Lefschetz, due to Mourougane \cite[Theorem~2.6 and Serre duality]{Mourougane1999}; we use the formulation \cite[Corollary~2.1.2]{DPS2001}. The determinant construction follows Lazi\'c--Peternell \cite[Lemma~4.1]{LP2018} and its later use in \cite[Lemma~5.1]{LMPTX2023}; we prove the unbounded-twist argument in the form needed here, including the numerically trivial case. The inverse determinant is pseudoeffective by the cotangent-tensor theorem of Lazi\'c--Matsumura--Peternell--Tsakanikas--Xie \cite[Theorem~4.1]{LMPTX2023}, which develops Ou's generic-nefness argument \cite[Theorem~1.4]{Ou2023}.

Berndtsson's smooth direct-image theorem \cite{Berndtsson2009} gives positivity of the fiberwise adjoint metrics. We obtain the continuous maximum metric by taking their normalized moment limits, and derive local horizontal strictness from his curvature formula. For singular coefficients, Berndtsson--P\u aun proved direct-image positivity over smooth families \cite[Theorem~3.5]{BP2008}; we use P\u aun--Takayama's multiplier-ideal formulation and extension to the actual torsion-free direct image \cite[Theorems~3.3.4--3.3.5]{PT2018}. Positivity of an invariant line is obtained from a specified holomorphic quotient; it is never inferred for an arbitrary singular subbundle. Fujino--Matsumura injectivity \cite{FM2016} and Nadel vanishing \cite{Nadel1990} supply the cohomological ending. Guan--Zhou strong openness \cite{GuanZhou2015} is used to keep the relevant ideal fixed when a strict metric is mixed with another one.

The harmonic argument belongs to a different regularity regime. Richberg smoothing \cite{Richberg1968} provides uniform approximations of continuous base weights, and Fujisawa's derived decomposition \cite[Theorem~1]{Fujisawa2015}, obtained from Takegoshi's harmonic representatives, applies to the smooth semipositive coefficient on $V$. Related vanishing theorems impose global curvature domination: Fujino \cite[Theorem~1.3]{Fujino2018} uses a smooth semipositive metric on the pullback of an ample base line, whereas Matsumura \cite[Theorem~1.3]{Matsumura2018} uses the pullback of a base K\"ahler form. Our harmonic proof starts instead from the kernel inclusion on a nonempty ordinary open set and treats the continuous base twist by uniform approximation. The reduced-flat construction uses weak semistable reduction \cite{AK2000}; its finite-group descent and normal-target norm are kept distinct from a finite-flat norm-line construction.

\subsection{Conventions and reading path}
We write $\sim$, $\sim_{\Q}$, and $\equiv$ for linear, rational linear, and numerical equivalence. A local frame of squared norm $e^{-\phi}$ has curvature $i\partial\bar\partial\phi$; factors of $2\pi$ are included when viewing curvature as a first Chern form. A continuous semipositive metric has continuous plurisubharmonic weights. Locally bounded weights have both upper and lower bounds and need not be continuous.

The action on sections is $(g\cdot s)(gx)=g(s(x))$. Algebraic-torus invariants agree with invariants under its maximal compact torus. Averaging smooth metric weights and K\"ahler forms over the compact torus preserves semipositivity. A rank assertion on a dense base-change open is not an assertion on singular fibers.

The shortest route to Corollary~\ref{thm:ordinary} is the transfer proof in Sections~\ref{sec:models}--\ref{inj:character-transfer}, its natural invariant application in Section~\ref{sec:cohomological-input}, and invariant-section descent in Section~\ref{sec:descent}. Section~\ref{inj:extremal-weights} also obtains exact source characters. The remaining constructions explain different ways to convert cohomology into sections. Section~\ref{inj:injectivity} gives an independent ordinary proof through singular injectivity; Sections~\ref{har:sec:harmonic}--\ref{har:sec:alteration} develop harmonic restriction and its geometric inputs; Sections~\ref{inj:localized} and~\ref{inj:stabilized} use exact multiplier ideals and actual direct-image lines. The final section also gives the stronger tensor-holonomy calculation.

The table compares the inputs and the step that completes each conversion. The source-character constructions in Sections~\ref{inj:extremal-weights} and~\ref{har:sec:extremal} additionally keep track of an extremal character.
\begin{center}
\small
\begin{tabular}{@{}>{\raggedright\arraybackslash}p{.25\textwidth}>{\raggedright\arraybackslash}p{.29\textwidth}>{\raggedright\arraybackslash}p{.37\textwidth}@{}}
\toprule
Input & Base construction & Conversion \\
\midrule
Unbounded shifted invariant sections; $-M$ pseudoeffective & Common section-ratio field & Fixed-ideal vanishing; Theorem~\ref{inj:transfer} \\
Natural invariant cohomology; nonzero Euler value & Fixed invariant system & Harmonic restriction; Section~\ref{har:sec:eulerroute} \\
Ordinary cohomology in one degree at unbounded powers & Curvature-null field and reduced flat alteration & Descent of a continuous line and its injection; Theorem~\ref{thm:alteration-conversion} \\
No positive structure-sheaf cohomology & Null-pole field or determinant semigroup & Maximum metric and harmonic restriction; Section~\ref{har:sec:fields} \\
No positive structure-sheaf cohomology & Null-pole field, expressed by degree zero & Local strictness for the exact direct-image ideal; Section~\ref{inj:localized} \\
Unbounded natural invariant twisted sections; $-M$ pseudoeffective & Toroidal invariant Iitaka model & Affine boundary orders and strong openness; Section~\ref{inj:stabilized} \\
\bottomrule
\end{tabular}
\end{center}
